% !TeX root = courtade-kumar.tex
% !TEX TS-program = pdfLaTeX
\section{Deferred proofs}

\subsection{The posterior form of the induction step}\label{app:closure}
All information quantities are measured in bits, and $\nu=4\mu(1-\mu)$ throughout. We write
\[
 \eta(t)=1-\kappa(t),\qquad p=(1-\rho)/2,
 \qquad q=1-\rho^2,\qquad
 \alpha=\frac{\rho^2-\kappa(\rho)}q\quad(0<\rho<1).
\]
Here $\eta(t)$ is the residual uncertainty of a bit with sign mean $t$;
this quantity is needed for the conditional entropy budgets below. The
profile for information lost through noise is
\[
 \mathcal B_\rho(\mu)
 =1-\kldivb{\mu}{1/2}-Q_\rho(\mu)
 =q\left[1-\kldivb{\mu}{1/2}+\alpha\nu\right].
\]
The divergence expansion \eqref{eq:app-entropyseries} implies $\alpha>0$ and
\[
 \eta(\rho)\nu\le\mathcal B_\rho(\mu)\le1-\kldivb{\mu}{1/2}.
\]
The endpoint profiles are $\mathcal B_1(\mu)=0$ and
$\mathcal B_0(\mu)=1-\kldivb{\mu}{1/2}$.

In \autoref{lem:closure}, let
$A_s=\Pr(B=1\mid S=s,W)$ and define
\[
\begin{aligned}
 J(a_0,a_1)&=1-\frac{\kldivb{a_0}{1/2}+\kldivb{a_1}{1/2}}2,\\
 M_p(a_0,a_1)&=J((1-p)a_0+pa_1,pa_0+(1-p)a_1).
\end{aligned}
\]
Complementing $B$ and exchanging the values of $S$ if needed, assume
\[
 \mu=\frac{\E A_0+\E A_1}{2}\ge\frac12,\qquad
 d=\E A_1-\E A_0\ge0,
 \qquad \mu+d/2\le1.
\]
Write
\[
\begin{aligned}
 J_* &=1-\frac{\kldivb{\mu-d/2}{1/2}+\kldivb{\mu+d/2}{1/2}}2,\\
 j&=\frac{\mathcal B_\rho(\mu-d/2)+\mathcal B_\rho(\mu+d/2)}2.
\end{aligned}
\]
The identities
\[
 \muti{B}{W\mid S}=J_*-\E J(A_0,A_1),\qquad
 \muti{B}{S_\rho,W}=1-\kldivb{\mu}{1/2}-\E M_p(A_0,A_1)
\]
show that \autoref{lem:closure} is equivalent to
\begin{equation}\label{eq:app-local}
 4d^2\le\nu,\quad \E J(A_0,A_1)\ge j
 \quad\Longrightarrow\quad
 \E M_p(A_0,A_1)\ge\mathcal B_\rho(\mu).
\end{equation}
Thus $J$ and $M_p$ represent conditional uncertainties before and after
recombining the sections. There is no independence assumption on $A_0,A_1$.

\subsection{Two local estimates}
Put
\[
 M_*:=1-\frac{\kldivb{\mu-\rho d/2}{1/2}+\kldivb{\mu+\rho d/2}{1/2}}2.
\]

\subsubsection*{Mean consistency}
\begin{lemma}\label{lem:app-consistency}
For $\mu<1$ and every finite joint posterior law with these means and $\E
J(A_0,A_1)\ge j$,
\begin{equation}\label{eq:app-Pbound}
 \E M_p(A_0,A_1)\ge
 P:=M_*-\left(1-\frac{qd^2}{a(\mu)+d^2}\right)(J_*-j).
\end{equation}
\end{lemma}
\begin{proof}
Apply \autoref{lem:consistency} to the posterior law. Since
\[
 \muti{B}{W\mid S}=J_*-\E J,\qquad
 \muti{B}{S_\rho}=1-\kldivb{\mu}{1/2}-M_*,
\]
it gives $\E M_p\ge M_*-(1-qd^2/(a(\mu)+d^2))(J_*-\E J)$.
The coefficient of $\E J$ is nonnegative, so $\E J\ge j$ proves the claim.
\end{proof}

\subsubsection*{The entropy budget and the gap}
\begin{lemma}\label{lem:app-budget}
If $d=|\E A_0-\E A_1|>0$ and $\E J(A_0,A_1)\ge j$, with $0\le j\le\eta(d)$,
let $\tau\in[d,1]$ satisfy $\eta(\tau)/\tau=j/d$. Then
\begin{equation}\label{eq:app-Cbound}
 \E M_p(A_0,A_1)\ge C:=\frac d\tau\eta(\rho\tau).
\end{equation}
\end{lemma}
\begin{proof}
We first prove the pointwise comparison. Define $\tau(e)$ by
$\eta(\tau(e))/\tau(e)=e$ and
$\phi(e)=\eta(\rho\tau(e))/\tau(e)$. The ratio $\eta(t)/t$ decreases strictly, since
\begin{equation}\label{eq:app-entropyratio}
 \left(\frac{\eta(t)}t\right)'=-\frac{\mathcal A(t)}{t^2},\qquad
 \mathcal A(t)=1-\frac{\ln(1-t^2)}{2\ln2}>0.
\end{equation}
For a pair with $t=|a_0-a_1|>0$ and $x=|a_0+a_1-1|$, put
$\mathcal H(x,t)=[\eta(x+t)+\eta(x-t)]/2$. We claim
\begin{equation}\label{eq:app-pointwise}
 \frac{M_p(a_0,a_1)}t\ge\phi\!\left(\frac{J(a_0,a_1)}t\right).
\end{equation}
In the interior $x>0$, $x+t<1$, let $u=-\mathcal H_x>0$ and $v=\mathcal
H-t\mathcal H_t>0$. The ratio $v/u$ decreases with $t$. Indeed, if
\[
 \Theta=v-\frac{1-x^2-t^2}{2x}u,
\]
direct differentiation gives $u_t>0$, $\Theta_t=tu/x$, and $(v/u)_t=-\Theta
u_t/u^2$. At zero gap,
\[
 \Theta(x,0)=\frac{2x\eta(x)-(1-x^2)\operatorname{atanh} x/\ln2}{2x}\ge0:
\]
the numerator vanishes at $0,1$ and has second derivative
$-2\operatorname{atanh} x/\ln2\le0$.

Now follow the contour $\mathcal H(x,t(x))/t(x)=e$ toward $x=0$. It exists all
the way: the ratio decreases in $t$ from infinity to $[1-\kldivb{x}{1/2}]/[2(1-x)]$, whose
derivative in $x$ is $-\ln x/[2(1-x)^2\ln2]>0$. Along this contour $t'=-tu/v$,
and
\[
 \frac{\mathrm d}{\mathrm d x}\frac{\mathcal H(x,\rho t(x))}{t(x)}
 =\frac1t\left[\frac{u(x,t)v(x,\rho t)}{v(x,t)}-u(x,\rho t)\right]\ge0.
\]
At $x=0$, $t=\tau(e)$, proving \eqref{eq:app-pointwise}. Boundary pairs follow
by continuity; $J=0$ gives the deterministic opposite pair directly.

Finally $\phi$ is increasing and convex. Its derivative is
\[
 \phi'(e)=\frac{\mathcal A(\rho\tau)}{\mathcal A(\tau)}\in[0,1].
\]
The ratio decreases with $\tau$: if $\mathcal A(\sqrt x)=\sum b_kx^k$ with
$b_0=1$, $b_k=1/(2k\ln2)>0$, the numerator of its derivative is
\[
 \sum_{i>k}(i-k)b_ib_k(\rho^{2i}-\rho^{2k})x^{i+k-1}\le0.
\]
As $\tau$ decreases with $e$, convexity follows. Since $\phi(e)\ge e$ and
$0\le\phi'(e)\le1$, a supporting line at $e=j/d$ has slope $\lambda\in[0,1]$
and intercept $c\ge0$, so \eqref{eq:app-pointwise} gives
\[
 M_p(a_0,a_1)\ge\lambda J(a_0,a_1)+c|a_0-a_1|.
\]
This also holds when $a_0=a_1$. Average and use $\E|A_0-A_1|\ge d$. The
resulting bound is $\lambda j+cd=C$. The zero-budget and noise endpoints
follow by continuity or directly.
\end{proof}

\subsection{Combining the estimates}
Assume $0<p<1/2$, $1/2\le\mu<1$, and $0<4d^2\le\nu$. Abbreviate
\[
 s=d^2,\quad a=a(\mu),\quad b_*=\mathcal B_\rho(\mu),
 \qquad j=qJ_*+q\alpha(\nu-s).
\]
One has $0<j\le J_*\le\eta(d)$, so both lemmas apply. 
\begin{proposition}[A compensated square]\label{lem:app-square}
With $C,P$ given by \eqref{eq:app-Cbound}, \eqref{eq:app-Pbound}, the positive weights
\[
 w_C=\frac{2(C+b_*-j)}q,\qquad
 w_P=\frac{a(1/\ln2+2\alpha \nu)}{\rho^2\nu}
\]
satisfy
\begin{equation}\label{eq:app-squarebound}
 w_C(C-b_*)+w_P(P-b_*)>qs(1-\nu)(\alpha-1/(8\ln2))^2\ge0.
\end{equation}
\end{proposition}
\begin{proof}[Proof of \autoref{lem:closure}]
Both weights are positive because $C\ge j$ and $b_*>0$. Thus $\E
M_p\ge(w_CC+w_PP)/(w_C+w_P)>b_*$, which proves \eqref{eq:app-local} for $d>0$.
If $d=0$, entropy concavity gives $\E M_p\ge\E J\ge b_*$. Deterministic means
are immediate. At $p=0$ the profile vanishes; at $p=1/2$ the budget and
Jensen's inequality force the sections to equal their means almost surely,
giving parent entropy $1-\kldivb{\mu}{1/2}$.
\end{proof}

\subsubsection*{The cancellation}
We prove \autoref{lem:app-square} using the scalar product lemma in
\autoref{app:scalar} and the elementary geometry in \autoref{app:geometry}.
Write
\begin{align*}
 \delta&=\frac{\muti{B}{S}}s
 =\frac{\kldivb{\mu-d/2}{\mu}+\kldivb{\mu+d/2}{\mu}}{2s},\\
 E_b&=1-\kldivb{\mu}{1/2}+a\delta-\frac{a+s}{2\nu\ln2},\\
 c_k&=\frac{\mu^{1-2k}+(1-\mu)^{1-2k}}{2^{2k+1}k(2k-1)\ln2}.
\end{align*}
The exact binary-information identity gives $s\delta=\sum_{k\ge1}c_ks^k$, with
$c_1=1/(2\nu\ln2)$ and $c_2=(4-3\nu)/(12\nu^3\ln2)$. The series converges also at $d=2(1-\mu)$.

Put $R=\rho^2$ locally. Expanding \eqref{eq:app-Pbound}, and retaining the
nonnegative remainder, gives
\begin{align}
 P-b_*&=\frac{Rqs}{a+s}E,\\
 E&\ge E_b-\frac\alpha R(a+q\nu+Rs)\nonumber\\
   &=E_b-\alpha(a+s)-\frac{q\alpha}R(a+\nu).
 \label{eq:app-Ebound}
\end{align}
The difference in the first inequality is exactly
\[
 R(a+s)\sum_{k\ge3}c_k(1+R+\cdots+R^{k-3})s^{k-1}\ge0.
\]

Define the scalar kernel
\begin{equation}\label{eq:app-kernel}
 K_p(t)=\frac{2\eta(\rho t)[\kappa(t)-\kappa(\rho t)]}{qt^2}.
\end{equation}
The identity $(C-b_*)(C+b_*-j)=C(C-j)-b_*(b_*-j)$ yields, for the left side
$\mathcal W$ of \eqref{eq:app-squarebound},
\begin{equation}\label{eq:app-exactblend}
 \frac{\mathcal W}s=K_p(\tau)-2q\left[1-\kldivb{\mu}{1/2}+\alpha\nu\right](\delta+\alpha)
 +\frac{qa(1/\ln2+2\alpha \nu)}{\nu(a+s)}E.
\end{equation}
Feasibility gives $\nu\ge d(2-d)$, and the chosen coordinate gives $\nu\ge4d^2$. Hence
\[
 \nu-s-\frac65d
 =\frac35[\nu-d(2-d)]+\frac25[\nu-4d^2]\ge0.
\]
Set $\beta=\min\{(\nu-s)/d,2\}\in[6/5,2]$. Since $\mathcal B_\rho(x)\ge
4x(1-x)\eta(\rho)$, the entropy ratio defining $\tau$ obeys
\[
 \frac{\eta(\tau)}\tau=\frac jd\ge\beta \eta(\rho).
\]
\autoref{lem:app-product} therefore bounds the first term of \eqref{eq:app-exactblend} by
$q[1/\ln2+3\alpha/\ln2+2(\beta+1)\alpha^2]$. \autoref{lem:app-channel} gives
\[
 \frac{q\alpha}R(1/\ln2+2\alpha \nu)\le \frac1{\ln2}-\frac1{2(\ln2)^2}.
\]
Insert this and \eqref{eq:app-Ebound} into \eqref{eq:app-exactblend}. The
multiplying factors are positive, so the inequality directions are preserved.
Collecting the three powers of $\alpha$ gives
\begin{equation}\label{eq:app-collect}
 \frac{\mathcal W}{qs}\ge F_0+F_1\alpha+2G\alpha^2,
\end{equation}
where, only for this calculation,
\begin{align*}
 F_0&=\frac1{\ln2}-2\left[1-\kldivb{\mu}{1/2}\right]\delta
 +\frac{a}{\nu(a+s)\ln2}\left[E_b-\left(1-\frac1{2\ln2}\right)(a+\nu)\right],\\
 F_1&=\frac3{\ln2}-2+2\kldivb{\mu}{1/2}-2\nu\delta+\frac{2a}{a+s}E_b-\frac a{\nu\ln2},\\
 G&=\beta+1-a-\nu.
\end{align*}
\autoref{app:geometry} proves the coupled bounds
\[
 F_0>\frac{1-\nu}{64(\ln2)^2},\qquad F_1> -\frac{1-\nu}{4\ln2},\qquad G\ge\frac{1-\nu}{2}.
\]
Consequently the right side of \eqref{eq:app-collect} is strictly larger than
\[
 (1-\nu)\left(\frac1{64(\ln2)^2}-\frac\alpha{4\ln2}+\alpha^2\right)
 =(1-\nu)(\alpha-1/(8\ln2))^2.
\]
The strictness holds also at $\mu=1/2$ because $F_0>0$. This proves
\autoref{lem:app-square}.

\section{A scalar entropy inequality}\label{app:scalar}
\begin{lemma}[Budget-constrained product]\label{lem:app-product}
Let $0<p<1/2$, $0<t<1$, and $6/5\le\beta\le2$. If $\eta(t)/t\ge\beta \eta(\rho)$, then
\begin{equation}\label{eq:app-productgoal}
 K_p(t)>q[1/\ln2+3\alpha/\ln2+2(\beta+1)\alpha^2].
\end{equation}
\end{lemma}
We first bound the channel parameters, then use a tangent at $t=\rho$ to
reduce the remaining inequality to two endpoints.
The logarithmic remainder used below is
\begin{equation}\label{eq:app-remainder}
 \mathcal R(y):=-\frac{1-y}{y}\ln(1-y)
 =1-\sum_{n\ge1}\frac{y^n}{n(n+1)}.
\end{equation}

\subsection{An entropy-square expansion}
Put $f(x)=\eta(\sqrt x)$ and $a_n=1/[2n(2n-1)\ln2]$. The divergence series
gives $f(x)=1-\sum_{n\ge1}a_nx^n$ and $\sum a_n=1$. Direct convolution shows
\begin{equation}\label{eq:app-squaretails}
 f(x)^2=1-x/\ln2+\sum_{n\ge2}b_nx^n,\qquad b_n>0,\qquad
 \sum_{k>n}b_k\ge a_n/\ln2\quad(n\ge3).
\end{equation}
Set $r_n=\sum_{j\ge n}a_j=\frac1{\ln2}\int_0^1x^{2n-2}/(1+x)\,\mathrm d x$ and
$O_{n-1}=\sum_{j=1}^{n-1}(2j-1)^{-1}$. Partial fractions give
\[
 b_n=2a_n\left(\frac{O_{n-1}}{(n-1)\ln2}-r_n\right)>0.
\]
Also $r_n\le1/[4(n-1)\ln2]$. For $n\ge5$,
\[
 \frac{(n-1)(\ln2)b_n}{a_n}\ge2O_{n-1}-\frac12
 \ge\frac{77}{30}>\frac{17}{7}
 \ge\frac{4n-3}{2n-3}
 =\frac{(n-1)(a_{n-1}-a_n)}{a_n}.
\]
The case $n=4$ follows from $\ln2<347/500$, giving $b_4-(a_3-a_4)/\ln2>1/[63000(\ln2)^2]$.
Summation proves the tail assertion. All endpoint series are absolutely
convergent.

To index the kernel by squared correlation, write
$\mathcal K_R(t)=K_{(1-\sqrt R)/2}(t)$.
\begin{lemma}\label{lem:app-channel}
For $0<R<1$, $q=1-R$, and $\alpha=f(R)/q-1$, one has $\alpha>0$, $\alpha' >0$, and
\begin{equation}\label{eq:app-channelbound}
 q\alpha(1/\ln2+2\alpha)\le\left(\frac1{\ln2}-\frac1{2(\ln2)^2}\right)R.
\end{equation}
For fixed $0<t\le1$, $\mathcal K_R(t)$ is concave in $R$.
\end{lemma}
\begin{proof}
Division of the entropy series gives $\alpha=\sum_{n\ge1}r_{n+1}R^n$, proving
positivity and monotonicity. Rewrite
\[
 q\alpha(1/\ln2+2\alpha)=(1/\ln2-4)f(R)+\frac{2f(R)^2}{1-R}
 +(1-R)(2-1/\ln2).
\]
Its constant coefficient is zero, its linear coefficient is $1/\ln2-1/[2(\ln2)^2]$, and its
quadratic coefficient is $[-1/12+2(\ln2-1/2)^2]/(\ln2)^2<0$. For $n\ge3$, the coefficient is
$(4-1/\ln2)a_n-2\sum_{k>n}b_k\le(4-3/\ln2)a_n<0$. This proves
\eqref{eq:app-channelbound}.

For concavity, fix $y=t^2$, $v=f(y)$. The function $f(x)^2-vf(x)$ vanishes at
$x=y$, and its coefficients of degree $n\ge2$ are $b_n+va_n>0$. Therefore
\[
 \mathcal K_R(t)=\frac{2(1-v)}y
 -\frac2y\sum_{j\ge1}R^j\sum_{n>j}(b_n+va_n)y^n.
\]
Every nonconstant coefficient is negative. Differentiation on compact
subintervals of $R<1$ proves concavity, also when $y=1$ by the convergent
tails.
\end{proof}
In particular, for $\beta\le2$ the target in \eqref{eq:app-productgoal} is at most
\begin{equation}\label{eq:app-targetcap}
 \frac q{\ln2}+3q\alpha(1/\ln2+2\alpha)
 \le\frac1{\ln2}-\left(\frac3{2(\ln2)^2}-\frac2{\ln2}\right)R
 <\frac1{\ln2}-\frac R{10(\ln2)^2}.
\end{equation}

\subsection{Noise at least one sixth}
For $p\ge1/6$, the budget implies $t<3/4$, since
\[
 \frac{\eta(3/4)}{3/4}\le\frac{3\ln2+1}{6\ln2}<\frac{31}{60\ln2}
 <\frac65\eta(2/3).
\]
The last inequality follows from \eqref{eq:app-remainder} and
$\ln6>7/4$.

Put $r=4/9$ and $x=t^2\le9/16$. The divergence series gives
\[
 \eta(\sqrt r\,t)\ge 1-\frac{rx}{2\ln2}-\frac{r^2x^2}{12(1-rx)\ln2},
\]
\[
 \frac{2[\kappa(t)-\kappa(\sqrt r\,t)]}{(1-r)t^2}
 \ge\frac1{\ln2}\left[1+\frac{1+r}{6}x+\frac{1+r+r^2}{15}x^2
       +\frac{1+r+r^2+r^3}{28}x^3\right].
\]
Both lower factors are positive. Multiply them, subtract $1/\ln2-r/[10(\ln2)^2]$, and clear
$(\ln2)^2(1-rx)$. Replacing $\ln2$ by $693/1000$ decreases the result and gives
\[
 \frac2{45}-\frac{12173}{162000}x+\frac{37079}{1215000}x^2
 +\frac{601643}{131220000}x^3-\frac{9815357}{413343000}x^4
 +\frac{6305}{1240029}x^5.
\]
Using $x^4\le(9/16)^2x^2$, this is at least
\[
 \frac2{45}-\frac{19}{250}x+\frac{x^2}{50}>0,
\]
as the last quadratic decreases and is positive at $x=9/16$. Thus $\mathcal
K_r(t)>1/\ln2-r/[10(\ln2)^2]$. Since $\mathcal K_0(t)=2\kappa(t)/t^2\ge1/\ln2$, concavity
supplies $\mathcal K_R(t)>1/\ln2-R/[10(\ln2)^2]$ on $0<R\le r$. Together with
\eqref{eq:app-targetcap}, this proves \eqref{eq:app-productgoal} for
$p\ge1/6$.

\subsection{The remaining noise range}
Put $b=1-p$, $e(t)=\eta(t)/t$, and $\zeta(t)=\operatorname{atanh}(bt)/(b\ln2)$.
Convexity of $\operatorname{atanh}$ gives
\begin{equation}\label{eq:app-midpoint}
 K_p(t)\ge K_{\rm mid}(t):=e(t)\zeta(t)+\frac q2\zeta(t)^2.
\end{equation}
We show that $K_{\rm mid}$ decreases on the entire interval $0<t<1$ when $b\ge5/6$.
Write $S(t)=\operatorname{atanh}(t)/(t\ln2)$ and $\mathcal A(t)=1-\ln(1-t^2)/(2\ln2)$.
Direct differentiation of $2bK_{\rm mid}$ gives
\[
 -\frac{2\operatorname{atanh}(bt)}{(\ln2)(1-b^2t^2)}
 \left[S(t)-b^2\mathcal A(t)-q/\ln2
 +\frac{\eta(t)}{t^2}\left(1-\frac1{(\ln2)S(bt)}\right)\right].
\]
Since $S(bt)\ge(1+b^2t^2/3)/\ln2$, the bracket is at least
$S-b^2\mathcal A-q/\ln2+b^2\eta/4$. As a quadratic in $b\in[5/6,1]$, this is a convex combination of
\begin{align*}
 U_0&=S-25\mathcal A/36+25\eta/144-5/(9\ln2),\\
 U_1&=S-5\mathcal A/6+5\eta/24-1/(3\ln2),\\
 U_2&=S-\mathcal A+\eta/4.
\end{align*}
All three are positive. The series gives $S-\mathcal A\ge0$, so $U_2>0$. In
$U_1$, only degrees one through three in $x=t^2$ can be negative; summing them
at $x=1$ and using $\ln2<347/500$ gives $U_1>1/(100\ln2)$. Its degree-$n$ coefficient
has numerator $16n^2-58n+15$, positive for $n\ge4$. In $U_0$, all coefficients
from degree two onward are positive (numerator $176n^2-338n+75$), and
\[
 U_0>\frac1{\ln2}\left[\frac2{25}-\frac{11}{100}x+\frac{x^2}{90}+\frac{x^3}{48}\right]
 \ge\frac7{3600\ln2}>0.
\]
The last polynomial decreases on $[0,1]$. Thus $K_{\rm mid}'<0$.

Let $t_\beta$ satisfy $e(t_\beta)=\beta H$, where $H=\eta(\rho)$. The budget gives
$t\le t_\beta$, and hence
\begin{equation}\label{eq:app-boundary}
 K_p(t)\ge K_{\rm mid}(t_\beta).
\end{equation}

\subsection{A tangent at the channel correlation}
As a function of $e$, $\zeta$ satisfies
\[
 \zeta_e=-\frac{t^2}{(\ln2)\mathcal A(t)(1-b^2t^2)},\qquad
 \zeta_{ee}=\frac{t^3[2(\ln2)\mathcal A(t)-t^2(1-b^2t^2)/(1-t^2)]}
 {(\ln2)^2\mathcal A(t)^3(1-b^2t^2)^2}.
\]
The second derivative is positive between $t_\beta$ and $\rho$. If
$\rho\ge5/6$, then $\beta\ge1/\rho$, so $e(t_\beta)=\beta H\ge e(\rho)$ and
$t_\beta\le\rho$, and on $t\le\rho$ its numerator bracket is at least
$2\ln2-t^4/16>0$, using $(1-b^2)/q\le9/16$ and $-\ln(1-t^2)\ge t^2+t^4/2$. If
$\rho<5/6$, then $H\ge\eta(5/6)$ and $\beta\ge6/5$, so $e(t_\beta)\ge e(5/6)$
and $t_\beta\le5/6$, and on this interval the bracket is at least
\[
 2(\ln2)\mathcal A(t)-\frac{t^2}{1-t^2}
 \ge\ln(144/11)-25/11>1/4.
\]
Consequently the tangent at $t=\rho$, where $e(\rho)=H/\rho$, gives
\begin{equation}\label{eq:app-canonicaltangent}
 \zeta(t_\beta)\ge\ell_\beta:=\frac{\operatorname{atanh}(b\rho)}{b\ln2}
 +HJ\left(\frac1\rho-\beta\right),\qquad
 J=\frac{\rho^2}{(\ln2)\mathcal A(\rho)(1-b^2\rho^2)}.
\end{equation}
\autoref{app:tangent} proves the endpoint estimates
\begin{equation}\label{eq:app-zbound}
 \ell_{6/5}-2\alpha>\frac{29}{50\ln2},\qquad
 \ell_2-2\alpha>\frac{1/4+3p/5+p^2/2}{\ln2}.
\end{equation}
In particular, $\ell_\beta>0$ throughout $6/5\le\beta\le2$, since it is affine
in $\beta$.

At $t=t_\beta$, write $\zeta=\zeta(t_\beta)$. Then
$K_{\rm mid}/q=\beta(1+\alpha)\zeta+\zeta^2/2$.
Inserting the positive tangent reduces its deficit to
\[
 Q_\beta=\beta(1+\alpha)\ell_\beta+\ell_\beta^2/2
          -1/\ln2-3\alpha/\ln2-2(\beta+1)\alpha^2.
\]
The coefficient of $\beta^2$ is $H^2J(qJ-2)/(2q)<0$: indeed $1-b^2\rho^2\ge q$
gives $qJ\le\rho^2/[(\ln2)\mathcal A(\rho)]\le1/\ln2<2$. It suffices to prove
$Q_{6/5},Q_2>0$.

Writing $z_\beta=\ell_\beta-2\alpha$ displays the cancellation:
\begin{equation}\label{eq:app-Qcancel}
 Q_\beta=\alpha[2\beta-3/\ln2+(\beta+2)z_\beta]
          +\beta z_\beta+z_\beta^2/2-1/\ln2.
\end{equation}
For $p\le1/6$, monotonicity of $\alpha$ and \eqref{eq:app-remainder} give
\[
 (\ln2)\alpha\ge\frac3{10}\left(\frac{43}{24}+\frac{41}{45}\right)
 -\frac{347}{500}>\frac19.
\]
Use \eqref{eq:app-zbound}, $693/1000<\ln2<347/500$, and \eqref{eq:app-Qcancel}.
At $\beta=6/5$ the result is
\[
 (\ln2)^2Q_{6/5}>\frac{7997}{562500}>0.
\]
For $\beta=2$, \autoref{app:tangent} proves positivity using the same lower
bound for $z_2$ and one elementary logarithm minorant. Thus both endpoint
deficits are positive, and concavity in $\beta$ proves
\eqref{eq:app-productgoal} for $p\le1/6$. This proves
\autoref{lem:app-product}.

\section{Polynomial coefficient bounds}\label{app:geometry}
We prove the three coupled coefficient bounds used in \eqref{eq:app-collect}.
In this appendix only, let
\[
 A=\frac a\nu=1+\mu-\frac{(1-\mu)^2}{3},\quad
 x=\frac s\nu,\quad c=1-\frac1{2\ln2},\quad E=1-\frac{1-\nu}{2\ln2}.
\]
The divergence series and $d\le2(1-\mu)$ give
\begin{equation}\label{eq:app-geomdomain}
 \frac{17}{12}\le A<2,\quad 1-\kldivb{\mu}{1/2}\le E,\quad
 c_2s\le \delta-c_1\le\frac c\nu,\quad
 x\le x_*:=\min\left\{\frac14,\frac{1-\mu}{\mu}\right\}.
\end{equation}
For the upper bound on $\delta-c_1$, use
$c_kd^{2k-2}\le[2k(2k-1)\nu\ln2]^{-1}$ and sum over $k\ge2$.
The two possible values of $x_*$ agree at $\mu=4/5$.

\subsection*{The constant coefficient}
Clearing the explicit logarithm denominators gives
\begin{align*}
 (A+x)(\ln2)^2F_0&=H_0+(\ln2)T(\delta-c_1),\\
 H_0&=(\ln2)(A+x)-c(\ln2)A(A+1)
       -\frac{x(\ln2)[1-\kldivb{\mu}{1/2}]}\nu-\frac{Ax}{2\nu},
\end{align*}
where $T=A^2-2(\ln2)[1-\kldivb{\mu}{1/2}](A+x)$. By \eqref{eq:app-geomdomain}, this is at least
the smaller of
\[
 H_+=H_0+T\frac{(4-3\nu)x}{12\nu^2},\qquad H_-=H_0+T\frac{c\ln2}\nu.
\]
Replacing $1-\kldivb{\mu}{1/2}$ by $E$ decreases both expressions; below, $H_\pm$ denote
these lower bounds. The first is concave in $x$, the second affine. At $x=0$,
$A>2\ln2\ge2(\ln2)[1-\kldivb{\mu}{1/2}]$, so their minimum divided by $A$ is $\ln2-c(\ln2)(A+1)>1/10>(1-\nu)/64$.
It remains to check $x=x_*$.

In fact $H_->(A+x)/64$ throughout the domain. Differentiate $\nu[H_--(A+x)/64]$
in $A$ and use $A\ge4/3$; the result is at least
\[
 -2(\ln2)^2+\ln2+31/64-x/2
 +(1-\nu)(11\ln2/3-445/192)>0.
\]
At $A=4/3$, the coefficient of $x$ is negative, so use $x=1/4$. The remaining
expression is
\[
 -19(\ln2)^2/6+19\ln2/12+365/768
 +(1-\nu)(28\ln2/9-5063/2304)>0,
\]
as follows by $0\le\nu\le1$ and $693/1000<\ln2<347/500$; a lower bound is $3289/562500$.

For $H_+$, subtract $(A+x)(1-\nu)/64$ and set $x=x_*$. Its derivative with respect to $\ln2$ is
$-A^2-x(1-\nu)/\nu-2(A+x)(4-3\nu)x/(12\nu^2)<0$.
Use $\ln2=347/500$. Clear the positive denominators in the table below, obtaining
the indicated polynomial $N$. The following elementary expansions prove
positivity:
\begin{center}\small
\begin{tabular}{@{}lll@{}}\hline
Range & Multiplier of $H_+-(A+x)(1-\nu)/64$ & Variable\\\hline
$1/2\le\mu\le4/5$ & $864000\nu^2$ & $u=8/5-2\mu\in[0,3/5]$\\
$4/5\le\mu<1$ & $27648000(1-\mu)\mu^4$ & $u=2\mu-8/5\in[0,2/5)$\\\hline
\end{tabular}
\end{center}
In the first row, expansion gives coefficient bounds
\begin{align*}
 N>&\ 900+6800u+943000u^2-1202500u^3
       -57500u^6-9500u^7-39u^8\\
  \ge&\ 900+6800u+200000u^2>0.
\end{align*}
The omitted fourth and fifth coefficients are positive. The last inequality
uses $u\le3/5$. In the second row the first six coefficients are positive, the
constant exceeds $5800$, and the final four terms give
\[
 N>5800+u^6(400000-14000u-10000u^2)>0.
\]
These bounds prove $F_0>(1-\nu)/[64(\ln2)^2]$.

\subsection*{The linear coefficient}
The coefficient of $\delta$ in
\[
 F_1=\frac3{\ln2}-\frac{2s[1-\kldivb{\mu}{1/2}]}{a+s}
       +2\left(\frac{a^2}{a+s}-\nu\right)\delta-\frac{2a}{\nu\ln2}
\]
is positive by $A>4/3$ and $x\le1/4$. Use $\delta\ge c_1+c_2s$, $1-\kldivb{\mu}{1/2}\le E$, and
multiply by $(A+x)\ln2$ to clear the logarithm denominator. Then
\begin{align*}
 (A+x)((\ln2)F_1+(1-\nu)/4)\ge J(x):={}&2(A+x)-A^2-2Ax-2x(\ln2)E\\
 &+\frac{(4-3\nu)x}{6\nu}(A^2-A-x)+\frac{1-\nu}{4}(A+x).
\end{align*}
This is concave in $x$, and $J(0)=A(2-A+(1-\nu)/4)>0$. At $x=x_*$ replace $\ln2$ by
$347/500$. On the two mean intervals, multiply respectively by
\[
 432000\nu,\qquad 864000\mu^3.
\]
For the first numerator, put $u=8/5-2\mu$. Its shifted coefficients imply
\[
 N>18000-117600u+(522000-470000u^2)u^2
 \ge352800(u-1/6)^2+8200>0.
\]
For the second, put $u=2\mu-8/5$. Its shifted coefficients imply
\[
 N>31000+u^3(660000-39000u-29000u^2-3000u^4)>0
 \quad(0\le u\le2/5).
\]
This proves $F_1>-(1-\nu)/(4\ln2)$.

\subsection*{The quadratic coefficient}
If the cap $\beta=2$ is active, then $A\le2-\nu/4$ gives
\[
 G=3-a-\nu\ge3(1-\nu)+\nu^2/4\ge(1-\nu)/2.
\]
Otherwise $G=(\nu-d^2)/d+1-a-\nu$ decreases with $d$. For $\mu\ge4/5$, substitute
the largest feasible gap $2(1-\mu)$ and use $a\le2\nu$ to get
\[
 G-(1-\nu)/2\ge\tfrac52(1-\nu)+4\mu-4\ge\tfrac1{10}>0.
\]
For $\mu\le4/5$, substitute $d=\sqrt \nu/2$ and use
$\sqrt\nu\ge(2+5\nu-\nu^2)/6$ (square both positive sides). Direct substitution of \eqref{eq:app-allowance} gives
\[
 G-(1-\nu)/2\ge\frac{12+(1-16\mu)\nu-4\nu^2}{12}.
\]
With $u=2\mu-5/4\in[-1/4,7/20]$, the numerator is
\[
 -4u^4+4u^3+\frac{39}{2}u^2+\frac34u+\frac3{64}
 \ge18(u+1/48)^2+\frac5{128}>0,
\]
since $u(1-u)\ge-5/16$. This completes the geometric square.

\section{Tangent endpoint estimates}\label{app:tangent}
This appendix proves \eqref{eq:app-zbound} and the remaining $Q_2$ sign on the
whole interval $0<p\le1/6$. We use the elementary logarithm bounds
\[
 \frac{693}{1000}<\ln2<\frac{347}{500},\quad
 \frac{43}{24}<\ln6<\frac95,\quad
 2<\ln12<\frac52,\quad
 \ln(3/2)<\frac{3041}{7500}<\frac{203}{500}.
\]
For example, the first four positive terms of
$\ln w=2\sum_{j\ge0}z^{2j+1}/(2j+1)$, $z=(w-1)/(w+1)$, with the remaining tail bounded by $2z^9/[9(1-z^2)]$, prove these after reducing $w$ by powers of two. Two terms with their tail give the displayed bound for $\ln(3/2)$.

For this appendix write $\ell=-\ln p$, $u=-\ln(1-p)$, $b=1-p$, $\rho=1-2p$, and
\[
 \Delta_p=(3-2p)(2-3p+2p^2),\qquad \mathcal R(p)=bu/p.
\]
At the canonical tangent, exact logarithm expansion gives
\begin{equation}\label{eq:app-canonicalexact}
\begin{aligned}
 Z:=\zeta(\rho)-2\alpha
 &=2+\frac{\ln[(2-3p+2p^2)/(3-2p)]-\mathcal R(p)}{2b\ln2},\\
 HJ&=\frac{2\rho^2}{\Delta_p\ln2}
       \left(1+\frac{\rho u}{p(\ell+u)}\right),
 \qquad z_\beta=Z+HJ(1/\rho-\beta).
\end{aligned}
\end{equation}

\subsection*{The first budget endpoint}
The logarithm ratio in \eqref{eq:app-canonicalexact} has derivative at least
$-6/7$, since its derivative is $(-5+12p-4p^2)/\Delta_p$ and
\[
 6\Delta_p+7(-5+12p-4p^2)=1+6p+44p^2-24p^3>0.
\]
With $\mathcal R(p)\le1-p/2-p^2/6$, it follows that
\[
 (\ln2)Z>Z_*(p):=\frac{1386}{1000}
 -\frac{1406/1000+5p/14-p^2/6}{2(1-p)}.
\]
Multiplication by $2(1-p)$ and a concave-quadratic endpoint check give
\[
 Z_*(p)>\begin{cases}
 17/25-p,&0<p\le1/12,\\
 7/10-6p/5,&1/12\le p\le1/6.
 \end{cases}
\]
For the first range, $u\le p+p^2/[2(1-p)]$ gives $\rho u/p<1$, while $\ell>2$.
Thus the parenthesis in \eqref{eq:app-canonicalexact} is less than $3/2$. The
prefactor $2\rho^2/\Delta_p$ decreases from $1/3$ to $3/14$ as $p$ increases to
$1/6$: factor it as $2[\rho/(3-2p)][\rho/(2-3p+2p^2)]$ and differentiate both
factors. Therefore $(\ln2)HJ<1/2$, and
\[
 (\ln2)z_{6/5}>\frac2{25}-p+\frac1{2\rho}\ge\frac{29}{50}.
\]
For the second range, convexity of $-\ln p$ gives $\ell\le16/5-42p/5$ from
its two endpoint bounds. Since $u\ge p+p^2/2$,
\[
 (3-7p)u/p-\ell\ge-\frac15+\frac{29}{10}p-\frac72p^2>0.
\]
Thus the parenthesis in \eqref{eq:app-canonicalexact} is at least $4/3$, and
$(\ln2)HJ\ge2/7$. It follows that
\[
 (\ln2)z_{6/5}>\frac35\rho+\frac2{7\rho}-\frac{17}{70}
 \ge2\sqrt{\frac6{35}}-\frac{17}{70}>\frac{29}{50}.
\]
All endpoint comparisons here are explicit quadratic or rational inequalities.

\subsection*{The second budget endpoint}
Define locally
\[
 l(p)=\frac{43}{24}+2\frac{1-6p}{1+6p},\qquad
 u_+(p)=p+\frac{p^2}{2(1-p)},\qquad z_*(p)=\frac14+\frac35p+\frac12p^2.
\]
The inequality $\ln x\ge2(x-1)/(x+1)$ for $x\ge1$ gives $\ell\ge l(p)$, and
$u\le u_+$. Put $a=p(3-2p)/2$, $c=2p/3$. Expanding the logarithm ratio gives
the lower bound
\[
 v(p):=-\frac{3041}{7500}-a-\frac{a^2}{2(1-a)}+c+\frac{c^2}{2}
 \le\ln\frac{2-3p+2p^2}{3-2p}.
\]
It follows from \eqref{eq:app-canonicalexact}, with every coefficient sign fixed, that
\[
 (\ln2)z_2\ge z_-(p):=
 \frac{1386}{1000}+\frac{v(p)-(1-p/2-p^2/6)}{2(1-p)}
 -\frac{2\rho(2\rho-1)}{\Delta_p}
       \left(1+\frac{\rho u_+}{p(l+u_+)}\right).
\]
The difference $z_--z_*$ equals $P_9(p)/\Delta_9(p)$, where
\begin{align*}
 \Delta_9={}&45000(1-p)(3-2p)(2-3p+2p^2)(91-97p+162p^2-72p^3)>0,\\
 P_9={}&295362-7540725p+84215700p^2-389321927p^3+673748366p^4\\
 &-511952624p^5+153218576p^6+16687440p^7-26964000p^8+6480000p^9.
\end{align*}
To prove its sign, put $v=1/6-p\in[0,1/6]$. Expansion and absorption of the
last two negative terms into $v^4$ give
\[
 P_9>32800+1000v\,[1290-20053v+67910v^2+311000v^3].
\]
Indeed the first four shifted coefficients exceed
$32800,1290000,-20053000,67910000$; the fourth-degree coefficient minus
$17244000/6^4+6480000/6^5$ exceeds $311000000$, and the intervening
coefficients are positive. For $w=v-1/10\ge-1/10$, the cubic in brackets is
\[
 \frac{1374}{5}+2859w+161210w^2+311000w^3
 \ge\frac{1374}{5}+2859w+130110w^2>250.
\]
The last inequality follows by completing the square; equivalently its
discriminant after subtracting $250$ is $-4733031$. Thus $(\ln2)z_2>z_*$.

It remains to prove $Q_2>0$. From \eqref{eq:app-remainder},
\[
 (\ln2)\alpha\ge a_-(p):=
 \frac{l(p)+1-p/2-p^2/[6(1-p)]}{4(1-p)}-\frac{347}{500}.
\]
The coefficient of $\alpha$ in \eqref{eq:app-Qcancel} is positive at
$\beta=2,z=z_*/\ln2$, and the expression increases with $z>0$. Hence
\begin{align*}
 (\ln2)^2Q_2
 &\ge a_-(4\cdot693/1000-3+4z_*)
       +2(693/1000)z_*+z_*^2/2-347/500\\
 &=\frac{N_7(p)}{3000000(1-p)^2(1+6p)},
\end{align*}
where
\begin{align*}
 N_7={}&218321-3964841p+20656844p^2-9548824p^3\\
 &-2992000p^4-12420000p^5+1275000p^6+2250000p^7.
\end{align*}
With $w=p-1/10$, $|w|\le1/10$, its shifted constant exceeds $18000$, its
linear coefficient has magnitude less than $140000$, and its quadratic
coefficient, after absorbing the higher powers, exceeds $16000000$. Explicitly
the shifted coefficients from degree two onward are
\[
 \frac{87454309}{5},\quad-11954249,\quad-8932000,\quad-11182500,
 \quad2850000,\quad2250000.
\]
Consequently
\[
 N_7>18000-140000|w|+16000000w^2
 =16000000(|w|-7/1600)^2+\frac{70775}{4}>0.
\]
Thus $Q_2>0$, completing the endpoint estimates.
